%% toc: P3 35 35.9
\section{Failures of compactness and the limits of finitary transfer}
\label{sec:p03-w03-p03-041-compactness-failures}

In first-order logic, compactness lets one pass from information about
finite pieces of a theory to a model of the whole theory. The same step
underlies several other first-order results: completeness with finite
proofs, the existence of infinite models for theories with arbitrarily
large finite models, the upward L\"owenheim--Skolem theorem, and the
impossibility of characterizing an infinite well-ordering. This section
locates precisely where that passage from finite to infinite breaks down
for \(L_{\omega_1\omega}\) (\cref{def:p02-lomega-one-omega}). It also
gives a criterion, due to Ebbinghaus, that measures failure of countable
compactness by the ordinals a logic can pin down.

\emph{Inputs.} The language \(L_{\omega_1\omega}\) with its satisfaction
clauses, first-order logic, and the ordinal notions of a well-ordering and
its order type. \emph{Conclusions.} A sublogic of \(L_{\omega_1\omega}\)
satisfying explicit closure conditions is countably compact exactly when no
countable set of its sentences pins down an infinite ordinal
(\cref{thm:p03-w03-p03-041-pinning-compactness}). Every countable ordinal
from \(2\) on is pinned down by one \(L_{\omega_1\omega}\)-sentence, so
countable compactness fails for \(L_{\omega_1\omega}\) for a structural
reason (\cref{thm:p03-w03-p03-041-countable-ordinals}). No single
\(L_{\omega_1\omega}\)-sentence characterizes the class of well-orders,
and \(w(L_{\omega_1\omega})=\omega_1\)
(\cref{thm:p03-w03-p03-041-undefinability} and
\cref{thm:p03-w03-p03-041-bounded}). Three further
first-order transfer principles fail for \(L_{\omega_1\omega}\)
(\cref{thm:p03-w03-p03-041-transfer}).

\subsection{What is already established}
\label{sec:p03-w03-p03-041-established}

Several of the basic failures are proved elsewhere in this handbook, and
this section builds on them rather than repeating them.
\begin{enumerate}
\item \emph{A countable counterexample to compactness.} The set
      \(\{\bigvee_{n<\omega}P_n\}\cup\{\neg P_n:n<\omega\}\) of
      sentences over nullary symbols is finitely satisfiable but has no
      model (\cref{thm:p02-lomega-one-omega-calibration}, verified in
      \cref{sec:p02-w02-p02-039-locations}). Barwise's countable version
      is part~3 of \cref{thm:p02-w02-p02-005-bw-extended-ls}. In this
      chapter, \cref{thm:p03-w03-p03-039-compactness-failure} gives a
      version with constant symbols. The failure of upward
      L\"owenheim--Skolem is proved there with three examples
      (\cref{thm:p03-w03-p03-039-upward-ls-failure}), and part~1 of
      \cref{thm:p03-w03-p03-039-ultraproducts} shows that an ultraproduct
      need not preserve an \(\mathcal L_{\omega_1,\omega}\)-sentence.
      Clauses~2 and~3 of \cref{thm:p03-w03-p03-041-transfer} below give
      further instances: an ultraproduct of finite structures in the empty
      signature that loses finiteness, and a sentence with a countably
      infinite model and no uncountable one.
\item \emph{The finitary-proof boundary.} No sound relation in which each
      consequence uses only finitely many premises is complete for
      countable theories of \(L_{\omega_1\omega}\)
      (\cref{thm:p02-w02c-p02-003-infinitary-rule}). The calculus
      \(\mathbf H_{\infty}\), whose rule (R3) has countably many premises,
      is complete for countable theories
      (\cref{thm:p02-w02c-p02-003-infinitary-completeness}). Each of its
      theorems follows from a countable part of the premises (part~2 of
      \cref{thm:p02-w02c-p02-003-hinf-soundness}). As a result, it misses
      an inconsistency that only an uncountable premise set produces
      (\cref{thm:p02-w02c-p02-003-hinf-uncountable}). Karp's completeness
      theorem for her own calculus is recorded with its source as
      \cref{thm:p02-karp-completeness}.
\item \emph{What survives.} Every countable satisfiable set of
      \(L_{\omega_1\omega}\)-sentences has a countable model (part~1 of
      \cref{thm:p02-w02-p02-005-bw-extended-ls}). In this chapter
      countable models are built from consistency properties instead of
      by compactness (\cref{thm:p03-w03-p03-039-extended-model-existence}),
      and for a countable signature and a countable fragment every
      structure has a countable substructure that is elementary for that
      fragment (\cref{thm:p03-w03-p03-039-downward-ls}). Interpolation and Beth
      definability hold (\cref{thm:p02-w02c-p02-003-interpolation} and
      \cref{thm:p02-w02c-p02-003-beth}).
\end{enumerate}
Ebbinghaus names Barwise compactness for countable admissible fragments
as the first substitute for compactness
\cite[Ch.~II, \S5.3, Example~1, p.~63]{Ebbinghaus1985ExtendedLogics}.
That theorem is not stated or proved in this section, and no result below
depends on it. The same holds for Hanf numbers and for the
L\"owenheim--Skolem bounds of the larger infinitary languages.

\subsection{Compactness properties and pinned ordinals}
\label{sec:p03-w03-p03-041-definitions}

\begin{hbconvention}{def:p03-w03-p03-041-setting}{Signatures, variables and orders in this section}
A signature \(\tau\) here is \emph{positive relational}: it has relation
symbols only, each of arity at least \(1\). Equality is logical.
Structures have nonempty domains. The variables of \(L_{\omega_1\omega}\)
are \(v_0,v_1,v_2,\ldots\). A single formula may bind all of them, since
it may be a countable conjunction. For a binary relation \(<\), the
\emph{field} of \(<^{\mathfrak A}\) is the set of elements \(a\) with
\(a<^{\mathfrak A}b\) or \(b<^{\mathfrak A}a\) for some \(b\). A
\emph{well-ordering of its field} is a strict linear order of the field in
which every nonempty subset of the field has a least element. Its order
type is then the ordinal isomorphic to the field under \(<^{\mathfrak A}\).
Ebbinghaus allows elimination of function symbols and then takes \(\tau\)
relational \cite[Ch.~II, p.~62]{Ebbinghaus1985ExtendedLogics}. Nullary
symbols are excluded here for the reason given in
\cref{def:p03-w03-p03-041-limits}.
\end{hbconvention}

\begin{hbdefinition}{def:p03-w03-p03-041-compactness}{Countable compactness and \(\kappa\)-compactness}
Let \(\mathcal F\) be a logic and \(\kappa\) an infinite cardinal.
\(\mathcal F\) is \emph{\(\kappa\)-compact} if for every signature
\(\tau\) and every set \(\Phi\) of \(\mathcal F\)-sentences of \(\tau\)
of cardinality at most \(\kappa\), the following holds. If each finite
subset of \(\Phi\) has a model, then \(\Phi\) has a model. It is
\emph{compact} if it is \(\kappa\)-compact for every infinite \(\kappa\)
\cite[Ch.~II, Definition~1.2.4(i)--(ii), p.~31]{Ebbinghaus1985ExtendedLogics}.
\(\aleph_0\)-compactness is Barwise's countable compactness property and
compactness is his strong compactness property
\cite[Ch.~I, \S1.4, p.~7]{Barwise1985BackgroundAims}. Ebbinghaus writes
\(\mathscr L\) and \(\mathscr L[\tau]\) for a logic and its
\(\tau\)-sentences. Here \(\mathcal F\) and \(\mathcal F[\tau]\) are used,
because \(\mathcal L\) denotes a signature elsewhere in the handbook.
\index{compactness!countable}\index{compactness!kappa@\(\kappa\)-}
\end{hbdefinition}

\begin{hbdefinition}{def:p03-w03-p03-041-pinning}{Pinned ordinals and well-ordering numbers}
Let \(<\in\tau\) be binary and \(\Phi\) a set of \(\mathcal F\)-sentences
of \(\tau\). Then \(\Phi\) \emph{pins down the ordinal \(\alpha\) via
\(<\)} if
\begin{enumerate}
\item[(i)] for every model \(\mathfrak A\) of \(\Phi\), \(<^{\mathfrak A}\)
      is a well-ordering of its field; and
\item[(ii)] some model \(\mathfrak A\) of \(\Phi\) has \(<^{\mathfrak A}\)
      a well-ordering of order type \(\alpha\).
\end{enumerate}
For a cardinal \(\kappa\geq1\), \(w_\kappa(\mathcal F)\) is the supremum
of the ordinals pinned down by sets of \(\mathcal F\)-sentences of
cardinality at most \(\kappa\), with \(w_\kappa(\mathcal F)=\infty\) when
these ordinals are unbounded. The \emph{well-ordering number} is
\(w(\mathcal F)=w_1(\mathcal F)\), which concerns single sentences
\cite[Ch.~II, Definition~5.2.1, pp.~60--61]{Ebbinghaus1985ExtendedLogics}.
In (ii) the order type is that of the field, as in (i).
\index{pinning down an ordinal}\index{well-ordering number}
\end{hbdefinition}

The ordinal \(1\) is never pinned down under this reading. A strict order
whose field is nonempty has at least two elements in its field. Every
\(n\geq2\) is pinned down by a single first-order sentence, and so is
\(0\). Hence \(w_\kappa(\mathcal F)\geq\omega\) whenever \(\mathcal F\)
contains first-order logic (\cref{thm:p03-w03-p03-041-pinning-compactness}).

\begin{hbdefinition}{def:p03-w03-p03-041-sublogic}{Sublogics of \(L_{\omega_1\omega}\), renaming and relativization}
The \emph{even renaming} \(\varphi^{e}\) of an
\(L_{\omega_1\omega}\)-formula \(\varphi\) replaces every occurrence of
each variable \(v_i\), free or bound, by \(v_{2i}\). No odd-indexed
variable occurs in \(\varphi^{e}\).

Let \(R\) be a binary symbol and \(\varphi\) a formula in which \(v_1\)
does not occur. The \emph{relativization} \(\varphi^{R}\) of
\(\varphi\) to the \(R\)-successors of \(v_1\) is defined by recursion on
the well-founded formula tree. Atomic formulas are unchanged. Negation,
every binary connective, \(\bigwedge\) and \(\bigvee\) commute with the
operation, so that, for instance,
\((\bigwedge\Psi)^{R}=\bigwedge\{\psi^{R}:\psi\in\Psi\}\). The
quantifier clauses are
\[
 (\exists y\,\psi)^{R}=\exists y\,\bigl(R(v_1,y)\land\psi^{R}\bigr),
 \qquad
 (\forall y\,\psi)^{R}=\forall y\,\bigl(R(v_1,y)\to\psi^{R}\bigr).
\]
Then \(\FV(\varphi^{R})\subseteq\FV(\varphi)\cup\{v_1\}\), and index sets
stay countable, so \(\varphi^{R}\) is an \(L_{\omega_1\omega}\)-formula.

A \emph{sublogic of \(L_{\omega_1\omega}\)} is an assignment
\(\mathcal F\) of a set \(\mathcal F[\tau]\) of
\(L_{\omega_1\omega}\)-formulas of \(\tau\) to each positive relational
signature \(\tau\). It must satisfy the following conditions.
\begin{enumerate}
\item[(F1)] \(\mathcal F[\tau]\) contains every first-order formula of
      \(\tau\), and \(\mathcal F[\tau]\subseteq\mathcal F[\tau']\) whenever
      \(\tau\subseteq\tau'\).
\item[(F2)] \(\mathcal F[\tau]\) is closed under \(\neg\), binary
      \(\land\), and \(\exists v\) for every variable \(v\).
\item[(F3)] If \(\varphi\in\mathcal F[\tau]\) is a sentence, then
      \(\varphi^{e}\in\mathcal F[\tau]\).
\item[(F4)] If \(\varphi\in\mathcal F[\tau]\) is a sentence in which
      \(v_1\) does not occur, and \(R\notin\tau\) is binary, then
      \(\varphi^{R}\in\mathcal F[\tau\cup\{R\}]\).
\end{enumerate}
Its \emph{sentences} are the members without free variables.
\(L_{\omega_1\omega}\) itself, with every formula included, is a
sublogic. First-order logic is a sublogic because renaming and
relativization turn first-order formulas into first-order formulas.
\end{hbdefinition}

The even renaming is used because a sentence of \(L_{\omega_1\omega}\)
may already use every variable. After renaming, \(v_1\) is free to serve
as the parameter of the relativization.

\begin{hblemma}{thm:p03-w03-p03-041-relativization}{proved=annex:p03-w03-p03-041-relativization}{Renaming, relativization and reducts}
Let \(\varphi\) be an \(L_{\omega_1\omega}\)-formula of \(\tau\).
\begin{enumerate}
\item For every \(\tau\)-structure \(\mathfrak A\) and assignment \(s\),
      \(\mathfrak A\entails\varphi^{e}[s]\) iff
      \(\mathfrak A\entails\varphi[s\circ e]\), where
      \((s\circ e)(v_i)=s(v_{2i})\). In particular \(\varphi\) and
      \(\varphi^{e}\) have the same models when \(\varphi\) is a sentence.
\item Suppose \(v_1\) does not occur in \(\varphi\) and \(R\notin\tau\).
      Let \(\mathfrak A\) be a structure for a signature containing
      \(\tau\cup\{R\}\) and let \(a\) be an element. Suppose the set
      \(F_a=\{b:(a,b)\in R^{\mathfrak A}\}\) is nonempty, and let
      \(\mathfrak A_a\) be the \(\tau\)-structure with domain \(F_a\)
      in which each symbol of \(\tau\) denotes its restriction to
      \(F_a\). For every assignment \(s\) with \(s(v_1)=a\) and
      \(s(v)\in F_a\) for each \(v\in\FV(\varphi)\),
      \[
       \mathfrak A\entails\varphi^{R}[s]
       \quad\Longleftrightarrow\quad
       \mathfrak A_a\entails\varphi[s'],
      \]
      for every assignment \(s'\) into \(F_a\) that agrees with \(s\) on
      \(\FV(\varphi)\).
\item If \(\tau\subseteq\tau'\) and \(\mathfrak A'\) is a
      \(\tau'\)-structure, then \(\mathfrak A'\entails\varphi[s]\) iff
      \(\mathfrak A'\restriction\tau\entails\varphi[s]\).
\end{enumerate}
\end{hblemma}

\subsection{Countable compactness as a bound on pinned ordinals}
\label{sec:p03-w03-p03-041-criterion}

The next theorem is Ebbinghaus's Proposition~5.2.4, which states that a
logic is \(\aleph_0\)-compact iff \(w_{\aleph_0}\) equals \(\omega\)
\cite[Ch.~II, Proposition~5.2.4, p.~62]{Ebbinghaus1985ExtendedLogics}.
It is given here for sublogics of \(L_{\omega_1\omega}\) with the closure
conditions (F1)--(F4). The regular-logic form, without an infinitary
syntax or a positive relational vocabulary restriction, is
\cref{thm:p03-general-pinning}. The equivalence with clause~2 below is
the observation that opens the
source's \S5.2 \cite[Ch.~II, p.~60]{Ebbinghaus1985ExtendedLogics}.

\begin{hbtheorem}{thm:p03-w03-p03-041-pinning-compactness}{proved=annex:p03-w03-p03-041-pinning-compactness}{Countable compactness and pinned ordinals}
Assume the Axiom of Choice, and let \(\mathcal F\) be a sublogic of
\(L_{\omega_1\omega}\). The following are equivalent.
\begin{enumerate}
\item \(\mathcal F\) is \(\aleph_0\)-compact.
\item No countable set of \(\mathcal F\)-sentences pins down an infinite
      ordinal.
\item \(w_{\aleph_0}(\mathcal F)=\omega\).
\end{enumerate}
In every case \(w_\kappa(\mathcal F)\geq\omega\) for every
\(\kappa\geq1\). If a set \(\Phi\) of \(\mathcal F\)-sentences pins down
an ordinal \(\alpha\geq\omega\), then \(\Phi\) together with one further
first-order sentence pins down \(\alpha+1\). If \(\mathcal F\) is not
\(\aleph_0\)-compact, then some countable set of \(\mathcal F\)-sentences
pins down \(\omega\). For first-order logic, \(w_\kappa=\omega\) for
every \(\kappa\geq1\).
\end{hbtheorem}

The proof turns a countable counterexample to compactness into a
well-ordering. Let \(\Phi=\{\varphi_n:n<\omega\}\) be finitely satisfiable
with no model. Introduce new binary symbols \(<\) and \(R\), and assign to
each point \(x\) of the field of \(<\) a fibre \(\{z:R(x,z)\}\). Require
the fibre of a point with at least \(n\) points at or below it to satisfy
\(\varphi_n\). A point with infinitely many predecessors would then carry
a model of all of \(\Phi\), so every point has finitely many. Conversely,
the \(k\)-th point can carry a model of
\(\{\varphi_0,\ldots,\varphi_{k+1}\}\). Ebbinghaus's set \(\Phi'\) has
these three clauses
\cite[Ch.~II, proof of Proposition~5.2.4, p.~62]{Ebbinghaus1985ExtendedLogics}.
In the other direction, a countably compact logic cannot pin down an
infinite ordinal, because new symbols naming a descending chain are
consistent with every finite part of the theory.

\subsection{The countable ordinals in \texorpdfstring{\(L_{\omega_1\omega}\)}{L omega1 omega}}
\label{sec:p03-w03-p03-041-ordinals}

\begin{hbdefinition}{def:p03-w03-p03-041-mu}{Ebbinghaus's ordinal formulas \(\mu_\beta\)}
In the signature \(\{<\}\), define for every countable ordinal \(\beta\)
\[
 \mu_\beta(x)=\forall y\Bigl(y<x\leftrightarrow
       \bigvee\{\mu_\gamma(y)\mid\gamma<\beta\}\Bigr)
\]
\cite[Ch.~II, Example~5.2.3(c), p.~62]{Ebbinghaus1985ExtendedLogics}.
The recursion is on \(\beta\). The variable \(y\) is free in
\(\mu_\gamma(y)\), so its bound variable is taken to be \(x\); two
variables therefore suffice. The empty disjunction is false, so
\(\mu_0(x)\) says that nothing lies below \(x\). For \(\beta<\omega_1\)
the index set is countable, so \(\mu_\beta\) is an
\(L_{\omega_1\omega}\)-formula with the single free variable shown. Let
\(\operatorname{LO}(<)\) be the first-order axioms for a strict linear
order of the whole domain, and for \(1\leq\alpha<\omega_1\) put
\[
 \sigma_\alpha=\operatorname{LO}(<)\land\forall x\bigvee\{\mu_\beta(x)\mid\beta<\alpha\}.
\]
This is the source's sentence. Its first conjunct is the source's phrase
``\(<\) is a linear ordering'', read as a strict order of the domain.
\end{hbdefinition}

\begin{hblemma}{thm:p03-w03-p03-041-mu}{proved=annex:p03-w03-p03-041-mu}{What \(\mu_\beta\) says in a linear order}
Let \(<\) be a strict linear order of the domain of \(\mathfrak A\). For
every countable \(\beta\) and every \(a\),
\(\mathfrak A\entails\mu_\beta[a]\) iff
\(\{b:b<a\}\) is well-ordered by \(<\) with order type exactly \(\beta\).
\end{hblemma}

By \cref{thm:p03-lindstrom-ordinal-code}, the formula \(\eta_\beta\) of
\cref{def:p03-lindstrom-ordinal} has the same meaning in strict linear
orders, so \(\mu_\beta\) and \(\eta_\beta\) are equivalent there. The
sentences differ. \(\operatorname{LO}(<)\land\eta^0_\alpha\) holds only in
orders isomorphic to \(\alpha\), while \(\sigma_\alpha\) holds in every
well-order of nonzero type at most \(\alpha\). For pinning down, the
weaker sentence suffices.

\begin{hbtheorem}{thm:p03-w03-p03-041-countable-ordinals}{proved=annex:p03-w03-p03-041-countable-ordinals}{\(L_{\omega_1\omega}\) pins down every countable ordinal from \(2\) on}
For \(1\leq\alpha<\omega_1\), the models of \(\sigma_\alpha\) are exactly
the structures \((A,<)\) that are well-orders of some type \(\delta\) with
\(1\leq\delta\leq\alpha\). For \(2\leq\alpha<\omega_1\), the sentence
\(\sigma_\alpha\) pins down \(\alpha\) via \(<\). Consequently
\(w(L_{\omega_1\omega})\geq\omega_1\), and \(L_{\omega_1\omega}\) is not
\(\aleph_0\)-compact.
\end{hbtheorem}

The last clause is a third proof of the failure of countable compactness.
It follows from \cref{thm:p03-w03-p03-041-pinning-compactness}, because
\(\sigma_\omega\) alone pins down an infinite ordinal. Ebbinghaus's
Example~5.2.3(c) states this lower bound together with the equality
\(w(L_{\omega_1\omega})=\omega_1\). He obtains the upper bound from
countable admissible fragments
\cite[Ch.~II, Example~5.2.3(c), pp.~61--62]{Ebbinghaus1985ExtendedLogics}.
The upper bound is proved below by another route
(\cref{thm:p03-w03-p03-041-bounded}). Example~5.2.3(c) says that
the sentence pins down every countable \(\alpha\neq0\). Under the field
reading of Definition~5.2.1, the case \(\alpha=1\) is excluded by the
remark after \cref{def:p03-w03-p03-041-pinning}. This case does not affect
the supremum.

\begin{hbexample}{def:p03-w03-p03-041-worked}{From a compactness failure to a well-ordering and back}
Let \(\tau=\{U_m:m<\omega\}\) with each \(U_m\) unary. Let
\(\varphi_0=\bigvee_{m<\omega}\exists x\,U_m(x)\) and
\(\varphi_{m+1}=\neg\exists x\,U_m(x)\). The set
\(\Phi=\{\varphi_n:n<\omega\}\) is the unary form of the counterexample in
\cref{thm:p02-lomega-one-omega-calibration}. A finite part mentions
\(\varphi_{m+1}\) only for \(m\) below some \(M\), and the one-point
structure with \(U_M\) true is a model of it. No model of \(\Phi\) exists,
because \(\varphi_0\) needs some \(U_m\) nonempty.

\emph{Forward.} Form \(\Phi'\) as in the proof of
\cref{thm:p03-w03-p03-041-pinning-compactness}. The point \(k\) of the
field must have a fibre satisfying \(\varphi_n\) for every \(n\leq k+1\).
Take the fibre of \(k\) to be the single point \(b_k\), with
\(U_{k+1}=\{b_k\}\) inside it and every other \(U_m\) empty there. This
fibre satisfies \(\varphi_0\) through the disjunct \(m=k+1\). It satisfies
\(\varphi_{m+1}\) for every \(m\leq k\), because each such \(U_m\) is
empty in the fibre. The structure with domain \(\omega\cup\{b_k:k<\omega\}\),
the usual order on \(\omega\), and \(R=\{(k,b_k):k<\omega\}\) is a model
of \(\Phi'\) of order type \(\omega\). Now suppose that in some model a
field point had infinitely many points below it. Its fibre would make every
\(U_m\) empty and also some \(U_m\) nonempty, which is impossible. So
\(\Phi'\) pins down \(\omega\).

\emph{Back.} Add unary symbols \(D_n\) with the first-order sentences
``\(D_n\) has exactly one element \(d_n\)'' and ``\(d_{n+1}<d_n\)''.
Every finite part of the enlarged set is satisfied in the model just built.
If the last index mentioned is \(N\), let \(D_n=\{N-\min(n,N)\}\). The
whole set has no model, because \(d_0>d_1>\cdots\) would be a descending
sequence in a well-ordered field. The result is a second countable
counterexample to compactness. In it, the infinitary information is
confined to the relativized clauses, and everything else is first-order.
\end{hbexample}

\subsection{The upper bound \texorpdfstring{\(w(L_{\omega_1\omega})\leq\omega_1\)}{w(L omega1 omega) <= omega1}}
\label{sec:p03-w03-p03-041-upper-bound}

Ebbinghaus obtains the upper bound \(w(L_{\omega_1\omega})\leq\omega_1\)
from an inequality for countable admissible fragments
\cite[Ch.~II, Example~5.2.3(c), p.~62]{Ebbinghaus1985ExtendedLogics}.
The bound is proved here by a different route, the undefinability of
well-ordering, which Marker proves with a consistency property and the
model existence theorem
\cite[\S4.4, Theorem~4.27 and Corollary~4.28, pp.~51--53]{Marker2014InfinitaryNotes}.
The presentation follows the January 2014 version of Marker's lecture
notes, not the later book. The consistency properties and the extended
model existence theorem used are those of
\cref{def:p03-w03-p03-039-consistency-property} and
\cref{thm:p03-w03-p03-039-extended-model-existence}. Their language
\(\mathcal L_{\omega_1,\omega}\) of \cref{def:p03-w03-p03-039-language}
has the same sentences as \(L_{\omega_1\omega}\). The positive relational
convention of \cref{def:p03-w03-p03-041-setting} does not apply in this
subsection. Signatures may contain function and constant symbols.

\begin{hbtheorem}{thm:p03-w03-p03-041-undefinability}{proved=annex:p03-w03-p03-041-undefinability}{Undefinability of well-ordering}
Assume the Axiom of Choice. Let \(\tau\) be a signature with a binary
relation symbol \(<\), and let \(\varphi\) be an
\(L_{\omega_1\omega}\)-sentence of \(\tau\). Suppose that for every
\(\alpha<\omega_1\) there are a model \(\mathfrak M\entails\varphi\) and a
map \(e\colon\alpha\to\dom\mathfrak M\) with
\(e(\xi)<^{\mathfrak M}e(\eta)\) whenever \(\xi<\eta<\alpha\). Then there
are a model \(\mathfrak N\entails\varphi\) with countable domain and a
map \(h\colon\mathbb Q\to\dom\mathfrak N\) with
\(h(q)<^{\mathfrak N}h(r)\) whenever \(q<r\) in \(\mathbb Q\)
\cite[Theorem~4.27, p.~51]{Marker2014InfinitaryNotes}.
\end{hbtheorem}

Marker states the hypothesis and conclusion as embeddings of
\((\alpha,<)\) and \((\mathbb Q,<)\) into \(<^{\mathfrak M}\) and
\(<^{\mathfrak N}\). The maps above are the form in which the proof uses
and produces them. The idea of the proof is to add constants
\(d_q\) for \(q\in\mathbb Q\) with the axioms \(d_q<d_r\) for \(q<r\),
and to call a finite set of further sentences acceptable when, for every
countable \(\alpha\), it can be satisfied in some model of \(\varphi\) by
placing the finitely many constants \(d_q\) it mentions on an increasing
chain at mutual distance at least \(\alpha\). Room for one more constant
between two placed ones is then always available at a larger \(\alpha\),
and the countable disjunctions are handled by the regularity of
\(\omega_1\).

\begin{hbcorollary}{thm:p03-w03-p03-041-bounded}{proved=annex:p03-w03-p03-041-bounded}{Single sentences pin down only countable ordinals}
Assume the Axiom of Choice.
\begin{enumerate}
\item Let \(\varphi\) be an \(L_{\omega_1\omega}\)-sentence of a signature
      with a binary symbol \(<\), and suppose that in every model
      \(\mathfrak M\) of \(\varphi\) the relation \(<^{\mathfrak M}\) is a
      well-ordering of its field. Then there is \(\alpha<\omega_1\) such
      that in every model of \(\varphi\) the field has order type less
      than \(\alpha\), and so at most \(\alpha\)
      \cite[Corollary~4.28, p.~52]{Marker2014InfinitaryNotes}.
\item \(w(L_{\omega_1\omega})=w_{\aleph_0}(L_{\omega_1\omega})=\omega_1\).
\end{enumerate}
\end{hbcorollary}

Part~2 is the equality of Ebbinghaus's Example~5.2.3(c)
\cite[Ch.~II, Example~5.2.3(c), pp.~61--62]{Ebbinghaus1985ExtendedLogics}.
Its lower half is \cref{thm:p03-w03-p03-041-countable-ordinals}. For
\(w_{\aleph_0}\) the step is that a countable set of sentences has the
same models as its conjunction, which is again an
\(L_{\omega_1\omega}\)-sentence. Each \(\sigma_\alpha\) of
\cref{def:p03-w03-p03-041-mu} pins down its own \(\alpha\). By part~1
the ordinals pinned down by one sentence all lie below one countable
ordinal, so the supremum \(\omega_1\) is not attained, and no sentence
pins down \(\omega_1\) or any larger ordinal. In the vocabulary
of \cref{def:p03-w03-p03-047-bounded}, \(L_{\omega_1\omega}\) is bounded,
by \(\omega_1\).

\begin{hbwarning}{def:p03-w03-p03-041-marker-gaps}{Completing the printed proof}
Marker leaves most of the verification of the consistency property as an
exercise \cite[Exercise~4.29, p.~53]{Marker2014InfinitaryNotes}, and the
annex supplies all of it. Three changes to the printed construction are
needed for this.
\begin{enumerate}
\item The printed collection consists of the sets
      \(\sigma_0\cup\{\varphi\}\cup\{d_q<d_r:q<r\}\) only
      \cite[proof after Corollary~4.28, p.~52]{Marker2014InfinitaryNotes}.
      Clause~C0 of the definition requires closure under subsets
      \cite[Definition~4.1, p.~39]{Marker2014InfinitaryNotes}, and a subset
      omitting \(\varphi\) is not of the printed form. The annex takes all
      subsets of such sets.
\item The printed condition~(*) bounds the gaps between the placed
      constants but asks for no room above the last one, and when no
      constant is placed it asks nothing of the model. Adding a constant
      \(d_r\) with \(r\) above every placed index then has no position
      available. The annex requires a well-ordered chain extending at
      least \(\alpha\) beyond the last placed constant.
\item The printed proof is the proof of Theorem~4.27. Corollary~4.28 is
      stated without a separate derivation, and the derivation is given
      in the annex.
\end{enumerate}
\end{hbwarning}

\subsection{Other finite-to-infinite transfers that fail}
\label{sec:p03-w03-p03-041-transfers}

Barwise observes that the upward half of the first-order
L\"owenheim--Skolem theorem is a compactness argument. He also notes that
few logics have a simple analogue of it
\cite[Ch.~I, \S2.6, p.~19]{Barwise1985BackgroundAims}. Compactness also
shows that finite and infinite cannot be separated in first-order logic
\cite[Ch.~I, \S1.2, p.~6]{Barwise1985BackgroundAims}. The next theorem
gives precise counterexamples for \(L_{\omega_1\omega}\).

\begin{hbtheorem}{thm:p03-w03-p03-041-transfer}{proved=annex:p03-w03-p03-041-transfer}{Three first-order transfers that fail for \(L_{\omega_1\omega}\)}
\begin{enumerate}
\item \emph{Arbitrarily large finite models.} In the empty signature, the
      sentence
      \[
       \varphi_{\mathrm{fin}}=\bigvee_{1\leq n<\omega}
          \exists v_1\cdots\exists v_n\,\forall y\bigvee_{i=1}^{n}y=v_i
      \]
      is true in exactly the finite structures. It has arbitrarily large
      finite models and no infinite model. A first-order theory with
      arbitrarily large finite models has an infinite model.
\item \emph{\L o\'s's theorem.} Assume the Axiom of Choice. Let
      \(\mathcal U\) be a nonprincipal ultrafilter on \(\omega\), and let
      \(\mathfrak M_n\) be a set of \(n+1\) elements in the empty
      signature. The ultraproduct \(\prod_{n<\omega}\mathfrak M_n/\mathcal U\)
      has as elements the classes of functions \(f\) with
      \(f(n)\in\mathfrak M_n\), where \(f\) and \(g\) are identified when
      \(\{n:f(n)=g(n)\}\in\mathcal U\). Every \(\mathfrak M_n\) satisfies
      \(\varphi_{\mathrm{fin}}\), but the ultraproduct does not. So the
      equivalence of \L o\'s's theorem, which holds for first-order
      sentences and drives \cref{thm:p03-compactness-ultraproduct}, fails
      for \(L_{\omega_1\omega}\)-sentences.
\item \emph{Upward L\"owenheim--Skolem.} In the signature of constants
      \(d_n\) (\(n<\omega\)), the sentence
      \(\forall x\bigvee_{n<\omega}x=d_n\) has a countably infinite model
      and no uncountable model. An infinite first-order structure
      \(\mathcal M\) has elementary extensions of every cardinality
      \(\kappa\geq\max(\abs{\Lang},\abs{\dom\mathcal M},\aleph_0)\)
      (\cref{thm:p02-upward-ls-elementary-extension}), so a first-order
      theory with an infinite model has an uncountable one.
\end{enumerate}
\end{hbtheorem}

Clause~3 uses constants, which the positive relational convention of
\cref{def:p03-w03-p03-041-setting} excludes. That convention governs only
\cref{thm:p03-w03-p03-041-pinning-compactness}. The same sentence with
\(\omega_1\) further constants, required to be pairwise distinct, is the
theory of \cref{thm:p02-w02c-p02-003-hinf-uncountable}. Every countable
part of it has a model, and the whole has none.
Clause~2 differs from part~1 of \cref{thm:p03-w03-p03-039-ultraproducts},
where every factor is the same countable structure and the lost sentence
is a covering disjunction over constants: here the factors are finite of
growing size and the lost sentence is the finiteness sentence itself.

\begin{hbwarning}{def:p03-w03-p03-041-limits}{What these results do and do not establish}
\begin{enumerate}
\item The syntactic version is
      \cref{thm:p03-w03-p03-041-pinning-compactness}, with (F1)--(F4).
      For regular logics, \cref{thm:p03-general-pinning} proves the
      general form of Ebbinghaus's Proposition~5.2.4. Its proof uses
      relativization and simultaneous substitution, not a prescribed
      infinitary formula grammar.
\item The earlier fibre construction uses positive arities: a nullary
      relation keeps the same truth value on every subdomain.
      The construction of \cref{thm:p03-general-fibres} instead gives
      each relation an index argument. A nullary \(R\) becomes unary
      \(R^\sharp(d)\), allowing different truth values in different fibres.
      Functions and constants are projected into the fibre. No separate
      axiom for each vocabulary symbol is needed.
\item Ebbinghaus states the upper bound
      \(w(L_{\omega_1\omega})\leq\omega_1\) as a consequence of the
      corresponding inequality for countable admissible fragments, which
      he asserts without proof at that place
      \cite[Ch.~II, Example~5.2.3(c), p.~62]{Ebbinghaus1985ExtendedLogics}.
      \Cref{thm:p03-w03-p03-041-bounded} proves the bound without
      admissible sets. The inequality for countable admissible fragments
      itself is not proved in this section.
\item The failure of countable compactness does not bring the
      L\"owenheim property down with it. Part~1 of
      \cref{thm:p02-w02-p02-005-bw-extended-ls} keeps it, and Barwise
      presents this pairing as an instance of Lindstr\"om's theorem
      \cite[Ch.~I, \S2.1, p.~14]{Barwise1985BackgroundAims}. In this
      handbook that theorem is \cref{thm:p03-lindstrom}, with its own
      regularity hypotheses. The present section does not derive its
      results from it.
\item Completeness theorems for \(L_{\omega_1\omega}\) exist, but for
      infinitary derivations
      (\cref{thm:p02-w02c-p02-003-infinitary-completeness},
      \cref{thm:p02-karp-completeness}). Barwise explains that a
      completeness theorem usually yields compactness because a
      derivation uses finitely many hypotheses
      \cite[Ch.~I, \S1.4, p.~7]{Barwise1985BackgroundAims}. This is the
      step that fails when derivations may use countably many premises.
\end{enumerate}
\end{hbwarning}

\emph{Relationships.} \Cref{thm:p03-w03-p03-041-pinning-compactness}
connects compactness to the well-ordering numbers of Ebbinghaus's \S5.2.
\Cref{thm:p03-w03-p03-041-mu} links his ordinal formulas to the order-type
test of \cref{thm:p03-lindstrom-ordinal-code}. The finitary-proof
boundary is \cref{thm:p02-w02c-p02-003-infinitary-rule}. Its
countable-premise counterpart is
\cref{thm:p02-w02c-p02-003-hinf-uncountable}. The first-order facts
contrasted in \cref{thm:p03-w03-p03-041-transfer} are
\cref{thm:p03-compactness-ultraproduct} and
\cref{thm:p02-upward-ls-elementary-extension}. Barwise reports that \(L_{\omega_1\omega}\)
has the interpolation property but not the Robinson consistency property,
and that in general countable compactness is much weaker than full
compactness \cite[Ch.~I, \S2.3, p.~16]{Barwise1985BackgroundAims}.
\Cref{thm:p03-w03-p03-047-bounded-omega} shows that a countably compact
regular logic has well-ordering number \(\omega\).
\Cref{thm:p03-w03-p03-041-bounded} places \(L_{\omega_1\omega}\) one
step above, with well-ordering number \(\omega_1\), and its proof replaces
compactness by the model existence theorem
\cref{thm:p03-w03-p03-039-extended-model-existence}. Well-ordering itself
is characterized by one sentence of \(\mathscr L_{\omega_1\omega_1}\)
(\cref{thm:p03-w03-p03-039-nadel-well-orders}), so the bound depends on
the quantifiers being finitary.

\emph{Exact sources.} Definitions~1.2.4(i)--(ii) and~5.2.1,
Example~5.2.3(c), Proposition~5.2.4 with its proof, and the substitutes of
\S5.3 are read from \textcite{Ebbinghaus1985ExtendedLogics}, Chapter~II
of the 2017 Cambridge reissue of \emph{Model-Theoretic Logics}, at
pp.~31, 60--63. Barwise's framing is read from
\textcite{Barwise1985BackgroundAims}, Chapter~I of the same volume, at
pp.~6, 7, 14, 16 and 19. Theorem~4.27, Corollary~4.28, the proof
after it and Exercise~4.29 are read from
\textcite{Marker2014InfinitaryNotes}, the author-hosted lecture notes
dated 10 January 2014, at pp.~51--53, with Definition~4.1 at p.~39. The
proofs in the annex are written for this
section. They use the satisfaction clauses of
\cref{def:p02-lomega-one-omega} and the cited first-order theorems.
